\chapter{回路図なしの修理}
% full version: chapters/17_debugging.tex

\pencilfig{17}{\pencilart{17}}{回路図のない基板：プローブが聞けるのは、何十億本もの配線のうちほんの数本だけ。}

動いている基板を回路図なしで渡され、修理してくれと頼まれたとしよう。エンジニアには4つの道具がある。信号を聞くためのプローブ。自分の信号を入れるための信号発生器。部品を外して何が壊れるかを見る手。そして制御レジスタへのアクセス。脳もまさにこれらの道具で調べられている。

\section*{ニューロン千個分のプローブ}

脳に刺した電極は、近くにある数個のニューロンのカチッを聞く。860億個に対して電極千本など無に等しいように思える。ところが、それで画面上のカーソルを動かすには十分なのだ。なぜか。隣り合うニューロンは似たノイズを出し、腕を動かす活動は実はわずか10から20個の共通変数で表せる。全員を聞く必要はなく、この数個の「声」を聞ければ足りる。

電極千本からの生の信号は毎秒数億ビットにもなり、頭蓋骨の下からの無線チャネルでは通しきれない。一方、カチッそのものが運ぶ情報はおよそ1000分の1だ。だからカチッをインプラント上で直接取り出し、それだけを外へ送るのが得策である。目と同じ圧縮の手法だ。

\section*{Neuralinkがしていること}

Neuralink社は問題の工学的な側面に取り組んだ。硬い針の代わりに細くしなやかな糸、それを埋め込むロボット、そして密閉されたワイヤレスのインプラントである。同社の公表論文に記載されているのはコンピュータとケーブルでつながったシステムで、圧縮と無線通信は計画とされている。両腕が麻痺した人がカーソルを操作したというデータは、今のところ主に同社自身の報告で知られている。それによればインプラントには数十本の糸に約1000のチャネルがあり、埋め込み後に一部の糸がずれた。アイデアそのものは新しくない。約100個の電極からなる硬いアレイを使う大学の研究グループは、すでに人がカーソルを操作し、思考の力で文字を「書き」、話そうとする試みを毎分約60語の速さで画面上の文章に変えることを可能にしている。

\section*{読み出しと書き込み}

デコーダは多数のニューロンのカチッの頻度を読み取り、毎秒10ビット未満の情報を得る。ニューロン自体が運ぶもののごく一部だ。その代わり、脳とデコーダは互いから学ぶ。人はデコーダに適応し、デコーダは人に適応する。互いに合わせ合う2人の学習者は、微妙な問題である。一人一人は安定でも、一緒になると揺れ出すことがある。

脳への書き込みは読み出しより難しい。電極からの電流は、狙った1個のニューロンではなく周りのすべてを興奮させる。光る画素のように光の点を「描く」ことはでき、人は簡単な文字を見分けられる。だが、脳が画像を見るように目の圧縮された符号を書き込むことは、まだ誰にもできない。

\section*{プローブに見えないもの}

電極が聞いているのは、データを運ぶ電話網だ。だが放送局、つまりドーパミン、セロトニン、ノルアドレナリン、アセチルコリンは物質の濃度であり、プローブには聞こえない。記憶のすべてが保存されている結合の強さも見えない。人が感じるような痛みも見えない。データは見えるがレジスタが見えないデバッガは、同じ入力がなぜ違う答えを生むのか、いつまでも首をかしげることになる。

\fullversion{17}
